\chapter{Discussion}
\label{chap:discussion}

% -------------------------------------------------------
\section{Value Beyond Performance Metrics}
% -------------------------------------------------------

The evaluation results demonstrate a performance advantage for the custom architecture
as case volume grows, but the headline metrics alone understate the practical value of
the design. Even in a hypothetical scenario where correctness and completeness scores
were identical across configurations, the custom architecture would remain preferable in
a production legal setting for reasons that no automatic metric captures.

\paragraph{The factsheet as a product artifact.}
The baseline agent's output is ephemeral: it produces a text response and leaves no
persistent structured representation of the case behind. The custom agent, by contrast,
produces and continuously maintains a structured factsheet as a natural side-effect of
normal operation. This factsheet—comprising parties, events, claims, damages, deadlines,
and document metadata—is a machine-readable, queryable representation of the case that
can be consumed directly by frontend systems without any additional LLM inference.
Concrete applications include timeline visualisations, party graphs, claims dashboards,
and deadline trackers, all of which can be rendered from the factsheet with no extra
cost.

\paragraph{Human oversight and auditability.}
The structured state enables a qualitatively different relationship between the lawyer
and the system. Rather than trusting a black-box text response, a practitioner can
inspect the factsheet directly, verify individual entries against source documents, and
correct errors at the element level. If the agent produces an incorrect answer, the
source of the error can often be traced to a specific factsheet entry—a mis-extracted
event date or an omitted party—and corrected there, improving all subsequent responses.
This kind of targeted, human-in-the-loop correction is not possible with a stateless RAG
agent, where errors are buried in the implicit context and must be addressed by
re-uploading or re-prompting.

\paragraph{Versioned case understanding.}
Because the factsheet is updated incrementally at each session boundary, it constitutes
a natural audit trail of how the agent's understanding of the case evolved over time.
Earlier versions can be retained and compared, providing a record of when specific
facts, parties, or claims entered the agent's knowledge base. In legal practice, where
the timing of knowledge can itself be legally significant, this version history has
potential evidentiary and procedural value.

\paragraph{Downstream system integration.}
The structured state lowers the integration barrier for adjacent legal tools and
workflows. A downstream system can query the factsheet via a standard database interface
to, for example, check for conflicts of interest against other open cases, trigger
calendar entries for extracted deadlines, or populate court filing templates with
party and claim data. These integrations require structured, typed data; they cannot be
built on top of free-form agent responses without a costly and fragile natural language
parsing layer.

\paragraph{Summary.}
The performance metrics reported in Chapter~\ref{chap:results} capture one dimension of
the architecture's value: how accurately the agent answers factual questions across a
progression of sessions. They do not capture the value of the structured case
representation itself—a persistent, inspectable, and reusable artifact that the custom
architecture produces as a consequence of its design. In production deployments with
real case volumes and real practitioner workflows, this latent value is likely to be
at least as significant as the raw performance advantage.

% -------------------------------------------------------
\section{Uncertainties and Confounding Factors}
% -------------------------------------------------------

The evaluation results support the conclusion that the custom agent outperforms the baseline on correctness and relevancy, but several sources of uncertainty affect the magnitude and generalisability of this finding. These are discussed below.

\paragraph{Question difficulty and scope.}
The evaluation queries were designed to be factual and document-grounded, but the level of specificity varies considerably across sessions. Some questions have unambiguous single-document answers, while others require synthesising information across multiple documents and timelines. Harder queries amplify the architectural advantage of the custom agent — the factsheet and structured retrieval tools are most useful precisely when the answer requires cross-document reasoning — while easier queries may show little to no difference between architectures. The observed aggregate performance gap therefore depends on the composition of the query set: a test set dominated by simple lookups would understate the custom agent's advantage, while one focused on complex synthesis would overstate it. The current evaluation mix likely sits somewhere in between.

\paragraph{Dataset size and complexity.}
Only two datasets are used in the evaluation, and their characteristics differ substantially. The results from TOSL-2024, the larger and more complex real case, show a more pronounced advantage for the custom agent than those from THRD-2021. This is consistent with the design hypothesis: the factsheet-based architecture is most beneficial when the document set is large and the case is complex, because those are the conditions under which naive retrieval is most likely to fail. With only two datasets, however, it is difficult to disentangle the effect of document volume from other case-specific characteristics. A larger benchmark spanning more cases and varying document volumes would provide stronger evidence for the scaling behaviour observed here.

\paragraph{Underlying language model capability.}
Figure~\ref{fig:deepeval-barplot-llms} shows that the performance gap between the custom and baseline agents narrows as the underlying LLM becomes more capable: for the most advanced model tested (Gemini 2.5 Pro), the two configurations are nearly indistinguishable. This suggests two non-exclusive interpretations. First, a sufficiently capable model may compensate for the structural disadvantage of the baseline by inferring missing context or tolerating retrieval noise more robustly. Second, the evaluation queries and datasets may not be sufficiently demanding to expose the limits of a powerful model: the questions can largely be answered from the retrieved chunks even without the accumulated factsheet context. Conversely, the performance gap is most pronounced for the smaller models, for which structured, pre-digested context provides the greatest marginal benefit. This implies that the value of the custom architecture is not fixed but model-dependent: it acts as a capability amplifier that matters most when the underlying model has limited reasoning capacity or context-handling ability.

\paragraph{Stochastic variability.}
LLM outputs are inherently stochastic. Each evaluation set was run five times to reduce the impact of this variability, and the non-parametric Mann-Whitney U test is robust to the non-normal, left-skewed score distributions observed. Nevertheless, the total number of unique query-session combinations is constrained by the cost and time of running large-scale multi-model evaluations, and the reported p-values should be interpreted with awareness of the limited sample size. The significance of correctness ($p = 8.5 \times 10^{-3}$) and relevancy ($p = 3.4 \times 10^{-4}$) results are robust, but the borderline result for completeness ($p = 8.3 \times 10^{-2}$) should be treated cautiously.